\chapter{Morphismes plats}\label{IV}
Nous donnons ici surtout les propri\'et\'es de platitude qui nous
ont servi dans les expos\'es pr\'ec\'edents. Une \'etude plus
d\'etaill\'ee se trouvera au Chapitre~IV des \og \'El\'ements de
G\'eom\'etrie Alg\'ebrique\fg en pr\'eparation\footnote{\Cf EGA IV 11 et 12.}, o\`u on \'etudie de fa\c con syst\'ematique
la situation suivante: $X$ \'etant localement de type fini sur~$Y$
localement noeth\'erien, et~$\cal{F}$ coh\'erent sur~$X$ et
$Y$-plat, donner des relations entre les propri\'et\'es de~$Y$,
celles de~$\cal{F}$, et celles des faisceaux coh\'erents induits
par~$\cal{F}$ sur les fibres de~$X\to Y$ (du point de vue notamment de
la dimension, de la dimension cohomologique, de la profondeur
etc.). On a notamment une fa\c con syst\'ematique d'obtenir
des th\'eor\`emes du type \emph{Seidenberg} ou \emph{Bertini}
(pour les sections hyperplanes). Le r\'esultat essentiel pour
l'application des m\'ethodes de platitude dans ce contexte est le
suivant (qui sera d\'emontr\'e plus bas): Si~$Y$ est int\`egre,
$X$ de type fini sur~$Y$, $\cal{F}$ coh\'erent sur~$X$, il existe un
ouvert non vide~$U$ de~$Y$ tel que~$\cal{F}$ soit $Y$-plat aux points
de~$X$ au-dessus de~$U$. Une deuxi\`eme fa\c con, sans doute
encore plus importante, dont la platitude s'introduit en
G\'eom\'etrie Alg\'ebrique, est la \emph{théorie de
descente}: voir par exemple les deux expos\'es de Grothendieck sur
le sujet au S\'eminaire Bourbaki\footnote{et, pour un expos\'e
plus d\'etaill\'e, les Expos\'es~\ref{VIII} et~\ref{IX} plus
bas.}. La platitude semble ainsi une des notions techniques centrales
en G\'eom\'etrie Alg\'ebrique.

Rappelons que la notion de platitude et fid\`ele platitude a
\'et\'e introduite par Serre dans GAGA. Un expos\'e des
\no \ref{IV.1}~et~\ref{IV.2} suivants se trouvera aussi dans Alg.\ Comm.\ de Bourbaki
(qui bien entendu, comme le titre du livre l'indique, ne se borne pas
au cas d'anneaux de base commutatifs)\footnote{N. Bourbaki,
Alg\`ebre Commutative, Chap.\, I (Modules Plats), Act.\ Sc.\ Ind 1290,
Paris, Hermann (1961).}.

Contrairement aux expos\'es pr\'ec\'edents, nous ne supposons
pas que les anneaux envisag\'es sont n\'ecessairement
noeth\'eriens.

\section{Sorites sur les modules plats}
\label{IV.1}

Un module $M$ sur l'anneau~$A$ est dit \emph{plat}
\index{plat, fid\`element plat (module)|hyperpage}(ou $A$-plat si on veut pr\'eciser~$A$)
si le foncteur
$$
T_M \colon N \mto M \otimes_{A}N
$$
(qui est en tous cas exact \`a droite) est
\emph{exact},
\ie transforme monomorphismes en monomorphismes. Il revient au
m\^eme de dire que le premier foncteur d\'eriv\'e \`a droite,
ou tous les foncteurs d\'eriv\'es \`a droite, sont nuls,
\ie que l'on a\enlargethispage{.5cm}
$$
\Tor^{A}_{1}(M,N)=0 \qquad \text{pour tout $N$,}
$$
\resp qu'on a
$$
\Tor^{A}_{i}(M,N)=0 \qquad \text{pour $i>0$, tout $N$.}
$$
Comme les $\Tor_{i}$ commutent aux limites inductives, il suffit
d'ailleurs de v\'erifier ces conditions pour $N$ de type fini, et
m\^eme (prenant alors une suite de composition de~$N$ \`a quotients
monog\`enes) que l'on ait
$$
\Tor^{A}_{1}(M,N)=0
$$
si $N$ monog\`ene, \ie de la forme $A/I$, $I$ \'etant un
id\'eal de~$A$. Notons d'ailleurs que
$$
\Tor^{A}_{1}(M,A/I)=0 \quad\Longleftrightarrow\quad I\otimes_{A}M \rightarrow
M=A\otimes_{A}M\text{ est \emph{injectif}},
$$
comme on voit sur la suite exacte des $\Tor$,
compte tenu de~$\Tor^{A}_{1}(M,A)=0$. Donc $M$ plat \'equivaut \`a
dire que pour tout id\'eal~$I$, l'homomorphisme naturel
$$
I\otimes_{A}M \to IM
$$
est un isomorphisme. Il suffit d'ailleurs de le v\'erifier pour $I$
de type fini, a fortiori il suffit de v\'erifier que le foncteur
$M\otimes$ est exact sur les modules \emph{de type fini}.

Comme chaque fois qu'on a un foncteur exact $T$, si on identifie, pour
un sous-objet $N'$ de~$N$, $T(N')$ \`a un sous-objet de~$T(N)$, on a
$$
T(N' \cap N'')=T(N') \cap T(N'')
$$
$$
T(N' + N'')=T(N') + T(N'')
$$
pour deux sous-objets $N'$, $N''$ de~$N$.

Une somme directe de modules plats, un facteur direct d'un module
plat, est plat. En particulier, $A$ \'etant plat, un module
\emph{libre}, donc aussi un module \emph{projectif}, est plat. Le
produit tensoriel de deux modules plats est plat, et si $M$ est plat
sur~$A$, alors $M\otimes_{A}B$ est plat sur~$B$ pour tout changement
de base
$A \to B$
(\`a cause de l'associativit\'e du produit tensoriel et du fait
qu'un compos\'e de foncteurs
exacts est exact). Si $M$ est plat sur~$B$, $B$ plat sur~$A$, alors
$M$ est plat sur~$A$ (m\^eme raison).

La suite exacte des $\Tor$, plus la \og commutativit\'e\fg du $\Tor$,
donne:

\begin{proposition}
\label{IV.1.1}
Soit $0 \to M' \to M \to M'' \to 0$ une suite exacte de~$A$-modules,
$M''$ \'etant plat. Alors
\begin{enumerate}
\item[(i)] cette suite reste exacte par tensorisation par n'importe
quel $A$-module $N$
\item[(ii)] pour que $M$ soit plat, il faut et il suffit que $M'$ le
soit.
\end{enumerate}
\end{proposition}

On peut donc dire que du point de vue du comportement par produits
tensoriels, les modules plats sont \og aussi bons\fg que les modules
libres ou projectifs (et la suite exacte de~\ref{IV.1.1} en
particulier est \og aussi bonne\fg que si elle splittait).

Soit $S$ une partie multiplicativement stable de~$A$, alors $S^{-1}A$
est plat sur~$A$, car $S^{-1}A \otimes N=S^{-1}N$ est un foncteur
exact en~$N$. Si $M$ est $A$-plat, alors $S^{-1}M=S^{-1}A \otimes M$
est $S^{-1}A$-plat, la r\'eciproque \'etant vraie si $M \to
S^{-1}M$ est un
isomorphisme,
\ie si
les $s \in S$ sont bijectifs dans $M$ (\`a cause de la
transitivit\'e de la platitude, $S^{-1}A$ \'etant plat sur~$A$).
Plus g\'en\'eralement, le cas d'un morphisme de pr\'esch\'emas
$X \to Y$ et d'un faisceau quasi-coh\'erent $\cal{F}$ sur~$X$ dont
on veut \'etudier la platitude par rapport \`a~$Y$ conduit \`a
la situation avec deux anneaux:

\begin{proposition}
\label{IV.1.2}
Soient $A \to B$ un homomorphisme d'anneaux, $M$ un $B$-module, $T$
une partie multiplicativement stable de~$B$.
\begin{enumerate}
\item[(i)] Si $M$ est $A$-plat, alors $T^{-1}M$ est $A$-plat (donc
aussi $S^{-1}A$-plat pour toute partie multiplicativement stable $S$
de~$A$ s'envoyant dans $T$).
\item[(ii)] Inversement, si $M_{\goth{n}}$ est plat sur~$A_{\goth{n}}$
pour tout id\'eal maximal $\goth{n}$ de~$B$, $M_{\goth{n}}$ est plat
sur~$A$ (ou, ce qui revient au m\^eme, sur~$A_{\goth{m}}$, o\`u
$\goth{m}$ est l'id\'eal premier de~$A$ image inverse de~$\goth{n}$)
alors $M$ est $A$-plat.
\end{enumerate}
\end{proposition}

On a en effet la formule, fonctorielle par rapport au $A$-module~$N$:
$$
T^{-1}M \otimes_{A} N=T^{-1}(M \otimes_{A} N)
$$
car les deux membres sont fonctoriellement isomorphes \`a $T^{-1}B
\otimes_{B} M \otimes_{B} N_{(B)}$ (avec $N_{(B)}=N \otimes_{A} B$) en
vertu des formules d'associativit\'e de~$\otimes$. Il s'ensuit
aussit\^ot que si $M \otimes_{A} N$ est exact en~$N$, il en est de
m\^eme de~$T^{-1}M \otimes_{A} N$ (comme compos\'e de deux
foncteurs exacts), d'o\`u~(i). Et il s'ensuit de m\^eme (ii), car
pour v\'erifier l'exactitude d'une suite de~$B$-modules, il suffit
de v\'erifier l'exactitude des localis\'es en tous les id\'eaux
maximaux de~$B$.

\begin{proposition}
\label{IV.1.3}
\begin{enumerate}
\item[(i)] Soit
$M$ un $A$-module plat. Si $x \in A$ est non diviseur
de~$0$ dans $A$, il est non-diviseur de~$0$ dans $M$. En particulier,
si $A$ est int\`egre, $M$ est sans torsion.
\item[(ii)] Supposons que $A$ soit int\`egre et que pour tout
id\'eal maximal $\goth{m}$ de~$A$, $A_{\goth{m}}$ soit principal
(par exemple $A$ anneau de Dedekind, ou m\^eme principal). Pour que
le $A$-module $M$ soit plat, il faut et
il
suffit qu'il soit sans torsion.
\end{enumerate}
\end{proposition}

On obtient (i) en notant que l'homoth\'etie $x$ dans $M$ s'obtient
en tensorisant par~$M$ l'homoth\'etie $x$ dans $A$. Pour (ii), on
peut supposer d\'ej\`a $A$ principal gr\^ace \`a~\ref{IV.1.2}
(ii); il faut montrer que si $M$ est sans torsion, alors pour tout
id\'eal $I$ de~$A$, l'injection $I \to A$ tensoris\'ee par $M$ est
une injection, ce qui signifie que le g\'en\'erateur $x$ de~$I$
est non diviseur de~$0$ dans~$M$, O.K.

\section{Modules fid\`element plats}
\label{IV.2}

Un foncteur $F$ d'une cat\'egorie dans une autre est dit
\emph{fidèle} si pour tout $X$, $Y$, l'application $\Hom (X,Y) \to
\Hom (F(X),F(Y))$ est injective. S'il s'agit d'un foncteur additif de
cat\'egories additives, il revient au m\^eme de dire que $F(u)=0$
implique $u=0$, et cela implique que $F(X)=0$ implique $X=0$. Pour que
$F$ soit \emph{fidèle et exact}, il faut et il suffit que la
condition suivante soit v\'erifi\'ee: pour toute suite $M' \to M
\to M''$ de morphismes dans $\cal{C}$, la suite transform\'ee
$F(M') \to F(M) \to F(M'')$
est exacte \emph{si} et \emph{seulement si} la
pr\'ec\'edente l'est. Ou encore: $F$ est exact, et $F(X)=0$
implique $X=0$. (N.B. pour pouvoir parler d'exactitude, il faut
supposer les cat\'egories en jeu \emph{abéliennes}). Supposons
qu'on ait une famille $(M_{i})$ d'objets non nuls de~$\cal{C}$ tels
que tout objet non nul de~$\cal{C}$ ait un sous-objet admettant un
quotient isomorphe \`a un $M_{i}$. Alors $F$ fid\`ele et exact
\'equivaut \`a: $F$ exact, et $F(M_{i}) \neq 0$ pour tout~$i$. Si
$\cal{C}$ est la cat\'egorie des modules sur un anneau~$A$, on peut
prendre par exemple pour $(M_{i})$ la famille des $A/\goth{m}$,
$\goth{m}$ parcourant les id\'eaux maximaux de~$A$. (En effet, tout
module non nul admet un sous-module non nul monog\`ene, donc
isomorphe \`a un $A/I$, $I$ id\'eal~$\neq A$, lequel par Krull
admet un quotient $A/\goth{m}$). De ces sorites, on d\'eduit en
particulier:

\begin{proposition}
\label{IV.2.1}
Soit $M$ un $A$-module. Conditions \'equivalentes:
\begin{enumerate}
\item[(i)] Le foncteur $M \otimes_{A}$ est fid\`ele et exact.
\item[(i~bis)]
$M$ est plat, et $M \otimes_{A} N=0$ implique $N=0$
\item[(i~ter)] $M$ est plat, et $M \otimes A/\goth{m} \neq 0$ pour
tout id\'eal maximal $\goth{m}$ de~$A$.
\item[(ii)] Pour toute suite d'homomorphismes $N' \to N \to N''$, la
suite tensoris\'ee par $M$ est exacte si et seulement si la suite
initiale l'est.
\end{enumerate}
\end{proposition}

On dit alors que $M$ est un $A$-module \emph{fidèlement plat}.
\index{fidelement plat (module)@fid\`element plat (module)|hyperpage}\index{plat, fid\`element plat (module)|hyperpage}En particulier, si $M$ est fid\`element plat, alors $N \to N'$ est
un monomorphisme (\'epimorphisme, isomorphisme) si et seulement si
l'homomorphisme tensoris\'e par~$M$ l'est. Un module fid\`element
plat est \emph{fidèle}, puisque l'homoth\'etie $f$ dans~$M$
s'obtient en tensorisant par~$M$ l'homoth\'etie $f$ dans~$A$.

On voit comme
dans~\ref{IV.1}
les propri\'et\'es de
transitivit\'e habituelles: le produit tensoriel de deux modules
fid\`element plats est fid\`element plat, si $M$ est
fid\`element plat sur~$A$, $M \otimes_{A}B$ est fid\`element plat
sur~$B$ pour toute extension de la base $A \to B$, si $B$ est une
$A$-alg\`ebre qui est fid\`element plate sur~$A$ et si $M$ est un
$B$-module fid\`element plat, c'est un $A$-module fid\`element
plat.

\begin{corollaire}
\label{IV.2.2}
Soit $A \to B$ un homomorphisme local d'anneaux locaux, $M$ un
$B$-module de type fini. Pour que $M$ soit fid\`element plat
sur~$A$, il faut et il suffit qu'il soit plat sur~$A$ et non nul.
\end{corollaire}

R\'esulte du crit\`ere (i~ter) et de Nakayama. En particulier,
\emph{pour que $B$ soit $A$-plat, il faut et il suffit qu'il soit
fidèlement $A$-plat}.

\begin{proposition}
\label{IV.2.3}
Soit $A \to B$ un homomorphisme d'anneaux, $M$ un $B$-module qui est
fid\`element plat sur~$A$. Pour tout id\'eal premier $\goth{p}$
de~$A$, il existe un id\'eal premier~$\goth{q}$ de~$B$ qui
l'induise.
\end{proposition}

Divisant par $\goth{p}$, on est ramen\'e au cas o\`u
$\goth{p}=0$. Localisant en l'id\'eal premier $0$, on est ramen\'e
au cas o\`u $A$ est un corps. Mais $M$ \'etant fid\`element plat
sur~$A$ est non nul, a fortiori $B \neq 0$, donc $B$ a un id\'eal
premier, qui ne peut qu'induire l'unique id\'eal premier de~$A$!
G\'eom\'etriquement, on peut dire que l'existence d'un faisceau
quasi-coh\'erent $\cal{F}$ sur~$X=\Spec (B)$ qui soit
\og fid\`element plat\fg relativement \`a~$A$, implique que $X \to
Y=\Spec (A)$ est \emph{surjectif}.

\begin{corollaire}
\label{IV.2.4}
Supposons $M$ plat sur~$A$, de type fini sur~$B$ et
$\Supp M=\Spec(B)$
(\ie $M_{\goth{q}}\neq 0$ pour tout id\'eal premier $\goth{q}$
de~$B$). Alors les id\'eaux premiers $\goth{q}$ de~$B$ contenant
$\goth{p}B$ minimaux induisent~$\goth{p}$.
\end{corollaire}

On
est encore ramen\'e au cas $\goth{p}=0$ (car les hypoth\`eses se
conservent toutes en divisant), donc $A$ int\`egre. On est
ramen\'e \`a l'\'enonc\'e suivant:

\begin{corollaire}
\label{IV.2.5}
$M$ \'etant comme dessus, tout id\'eal premier minimal $\goth{q}$
de~$B$ induit un id\'eal premier $\goth{p}$ de~$A$ qui est minimal.
\end{corollaire}

En effet, localisant en $\goth{p}$ et $\goth{q}$, on est ramen\'e
\`a prouver que si $A$ et $B$ sont locaux et l'homomorphisme $A \to
B$ local, $M$ un $B$-module non nul plat sur~$A$, et si $B$ est de
dimension $0$, alors $A$ est de dimension $0$. Par~\ref{IV.2.2}
et~\ref{IV.2.3}, on conclut que tout id\'eal premier de~$A$ est
induit par un id\'eal premier de~$B$, donc par l'id\'eal maximal
de~$B$,
donc est l'id\'eal maximal, cqfd. G\'eom\'etriquement,
\ref{IV.2.5} signifie que toute composante irr\'eductible de
$X=\Spec (B)$
domine quelque composante irr\'eductible de~$Y=\Spec (A)$ (moyennant
l'existence d'un faisceau quasi-coh\'erent de type fini sur~$X$, de
support~$X$, et plat par rapport \`a~$Y$).

On notera qu'on n'a pas eu \`a supposer dans~\ref{IV.2.4} $M$
fid\`element plat sur~$A$, mais rien ne garantit alors l'existence
d'un id\'eal premier contenant $\goth{p}B$ (donc d'un minimal parmi
de tels).

\begin{proposition}
\label{IV.2.6}
Soit $i \colon A \to B$ un homomorphisme d'anneaux. Conditions
\'equivalentes:
\begin{enumerate}\item[(i)] $B$ est un $A$-module fid\`element plat.
\item[(ii)] $B$ est plat sur~$A$, et $\Spec (B) \to \Spec (A)$ est
surjectif
\item[(ii~bis)] $B$ est plat sur~$A$, et tout id\'eal maximal est
induit par un id\'eal de~$B$.
\item[(iii)] $i$ est injectif et $\Coker i$ est un $A$-module plat.
\item[(iv)] Le foncteur $M_{(B)}=M \otimes_{A} B$ en le $A$-module $M$
est exact, et l'homomorphisme fonctoriel canonique $M \to M_{(B)}$ est
injectif.
\item[(iv~bis)] Pour tout id\'eal $I$ de~$A$, $I \otimes_{A} B \to
IB$ est un isomorphisme, et l'image inverse de~$IB$ dans $A$ est
\'egale \`a~$I$.
\end{enumerate}
\end{proposition}

On a (i)$\To$(ii) par~\ref{IV.2.3},
(ii)$\To$(ii~bis) est trivial, (ii~bis)$\To$(i)
par le crit\`ere (i~ter) de~\ref{IV.2.1}. On a
(iii)$\To$(iv) par~\ref{IV.1.1}, (iv)$\To$(iv~bis)
trivialement (en faisant $M=A/I$ dans la deuxi\`eme
condition~(iv~bis)), et (iv~bis)$\To$(i) en vertu du
crit\`ere de platitude par id\'eaux vu au d\'ebut de~\ref{IV.1}
et du crit\`ere \ref{IV.2.1}~(i~ter). Enfin,
(iv)$\To$(iii) par une r\'eciproque facile
de~\ref{IV.1.1}, et (i)$\To$(iv) car
si $N$ est le noyau de~$M \to M \otimes_{A}B =T(M)$, alors ($T$~\'etant exact) $N \to T(N)$ est nul d'o\`u $T(N)=N \otimes_{A}B
=0$, d'o\`u $N=0$, cqfd.

\section{Relations avec la compl\'etion}
\label{IV.3}
Soient $A$ un anneau noeth\'erien, $I$ un id\'eal dans $A$,
$\widehat{A}$ le s\'epar\'e compl\'et\'e de~$A$ pour la
topologie $I$-pr\'eadique, et pour tout $A$-module $M$, soit
$\widehat{M}$ son compl\'et\'e pour la topologie
$I$-pr\'eadique. C'est un $\widehat{A}$-module, d'o\`u un
homomorphisme canonique
$$
M \otimes_A \widehat{A} \to \widehat{M}.
$$
Lorsque $M$ parcourt les modules \emph{de type fini}, le foncteur $M
\mto \widehat{M}$ est exact, comme il r\'esulte facilement du
\emph{théorème de Krull: Si $N \subset M$, la topologie de~$N$
est celle induite par la topologie de~$M$}. Comme $M \otimes_A
\widehat{A}$ est exact \`a droite, on en conclut ais\'ement (en
r\'esolvant $M$ par $L \to L' \to M$, avec $L$ et $L'$ libres de
type fini) que l'homomorphisme fonctoriel plus haut est un
\emph{isomorphisme} ($\widehat{M}$ \'etant aussi exact \`a droite)
et par cons\'equent que $M \otimes_A \widehat{A}$ est aussi un
foncteur \emph{exact} en $M$. Par suite:
\begin{proposition}
\label{IV.3.1}
Soient $A$ un anneau noeth\'erien, $I$ un id\'eal de~$A$, alors le
compl\'et\'e s\'epar\'e $\widehat{A}$ de~$A$ (pour la
topologie $I$-pr\'eadique) est \emph{plat} sur~$A$.
\end{proposition}

\begin{corollaire}
\label{IV.3.2}
Pour que $\widehat{A}$ soit fid\`element plat sur~$A$, il faut et il
suffit que $I$ soit contenu dans le radical de~$A$.
\end{corollaire}

En effet, il suffit d'appliquer le crit\`ere \ref{IV.2.1} (i~ter).

Ces r\'esultats r\'esument tout ce qu'on sait dire, du point de
vue de l'alg\`ebre lin\'eaire, sur les relations entre $A$ et
$\widehat{A}$. Le corollaire \ref{IV.3.2} est surtout utilis\'e lorsque $A$
est un anneau local noeth\'erien et que $I$ est contenu dans
l'id\'eal maximal $\goth{m}$ (et le plus souvent, lui est \'egal).

\section{Relations avec les modules libres}
\label{IV.4}

\begin{proposition}
\label{IV.4.1}
Soient $A$ un anneau, $I$ un id\'eal de~$A$, $M$ un
$A$-module. Supposons qu'on soit sous l'une ou l'autre des
hypoth\`eses suivantes:
\begin{enumerate}
\item[(a)] $I$ est nilpotent
\item[(b)] $A$ est noeth\'erien, $I$ est dans le radical de~$A$, et
$M$ est de type fini.
\end{enumerate}
Pour que $M$ soit libre sur~$A$, il faut et il suffit que $M \otimes
A/I$ soit libre sur~$A/I$, et que $\Tor_1^A(M,A/I) = 0$.
\end{proposition}

C'est
n\'ecessaire, prouvons la suffisance. Soit $(e_i)$ une famille
d'\'el\'ements de~$M$ dont l'image dans $M \otimes A/I = M/IM$ y
d\'efinit une base sur~$A/I$ (c'est une famille finie dans le cas
(b)). Soit $L$ le $A$-module libre construit sur le m\^eme ensemble
d'indices, on a donc un homomorphisme $L \to M$ tel que la
tensorisation $T$ par $A/I$ induit un isomorphisme $T(L)
\isomto T(M)$. Si $Q$ est le conoyau de~$L \to M$, on a
donc $T(Q) = 0$, d'o\`u $Q=0$ en vertu de Nakayama (valable sous
l'une ou l'autre condition (a) ou (b)). Donc $L \to M$ est surjectif,
soit $R$ son noyau, on a donc une suite exacte
$$
0 \to R \to L \to M \to 0
$$
d'o\`u, comme $\Tor_1^A(M,A/I)=0$, une suite exacte $0 \to T(R) \to
T(L) \to T(M) \to 0$, d'o\`u $T(R)=0$, d'o\`u encore
$R=0$
en vertu de Nakayama (tenant compte que dans le cas (b), $R$ est de
type fini puisque $A$ \'etait suppos\'e noeth\'erien).

\begin{corollaire}\label{IV.4.2}
On peut remplacer la condition $\Tor_1^A(M,A/I)=0$ par:
l'homomorphisme canonique surjectif
\begin{equation*}
\label{eq:IV.2.*}
\tag{$*$} { \gr_I^0(M) \otimes_{A/I} \gr_I(A) \to \gr_I(M) }
\end{equation*}
est un isomorphisme.
\end{corollaire}

En effet, si $M$ est libre, cela est certainement v\'erifi\'e. Il
faut donc prouver que si $M \otimes A/I$ est libre sur~$A/I$ et la
condition sur les $\gr$ v\'erifi\'ee, alors $M$ est libre. On
reprend la d\'emonstration ci-dessus en construisant $L \to M$, il
r\'esulte de l'hypoth\`ese que cet homomorphisme induit un
isomorphisme pour les gradu\'es associ\'es, donc son noyau est
contenu dans l'intersection des $I^nL$, donc est nul (comme il est
trivial dans (a), et bien connu dans (b)). Cqfd.



\begin{corollaire}
\label{IV.4.3}
Supposons que $A/I$ soit un corps. Alors les conditions suivantes sur~$M$ sont \'equivalentes:
\begin{enumerate}
\item[(i)] $M$ est libre
\item[(ii)] $M$ est projectif
\item[(iii)] $M$ est plat
\item[(iv)] $\Tor_1^A(M, A/I) = 0$
\item[(v)] L'homomorphisme canonique \eqref{eq:IV.2.*} est bijectif.
\end{enumerate}
\end{corollaire}

En effet, dans le cas envisag\'e, $M \otimes A/I$ est
automatiquement libre.

Le r\'esultat pr\'ec\'edent est valable dans les deux cas
suivants:
\begin{enumerate}
\item[(a)] $M$ est un module \emph{quelconque} sur un anneau local $A$
dont l'id\'eal maximal $I$ est \emph{nilpotent} (par exemple un
anneau local artinien).
\item[(b)] $M$ est un module \emph{de type fini} sur un anneau
\emph{local noethérien}.
\end{enumerate}

Rappelons
pour m\'emoire:
\begin{corollaire}
\label{IV.4.4}
Supposons que $A$ soit un anneau \emph{local noethérien
intègre} d'id\'eal maximal $\goth{m} = I$, de corps r\'esiduel
$k = A/I$, de corps des fractions $K$. Soit $M$ un module de type fini
sur~$A$. Alors les conditions \'equivalentes \textup{(i)} \`a \textup{(v)}
pr\'ec\'edentes \'equivalent aussi \`a
\enlargethispage{.5cm}\begin{enumerate}
\item[(vi)] $M \otimes_A K$ et $M \otimes_A k$ sont des espaces
vectoriels de m\^eme dimension (\ie le rang de~$M$ sur~$A$ est
\'egal au nombre minimum de g\'en\'erateurs du $A$-module $M$).
\end{enumerate}
\end{corollaire}

D\'emonstration imm\'ediate; on laisse au lecteur le soin de
g\'en\'eraliser au cas o\`u~$A$ est seulement suppos\'e sans
\'el\'ements nilpotents: il faut alors exiger que les rangs de~$M$
pour les id\'eaux premiers minimaux de~$A$ soient \'egaux \`a la
dimension de l'espace vectoriel $M \otimes_A k$.

\section{Crit\`eres locaux de platitude}
\label{IV.5}

\begin{proposition}
\label{IV.5.1}
Soit $A$ un anneau muni d'un id\'eal $I$, $M$ un
$A$-module. Supposons
$$
\Tor_1^A(M, A/I^n) = 0 \qquad \text{pour $n>0$}
$$
alors l'homomorphisme canonique surjectif
\begin{equation*}
\label{eq:IV.5.1.*}
\tag{$*$} {\gr_I^0(M) \otimes_{A/I} \gr_I(A) \to \gr_I(M)}
\end{equation*}
est un isomorphisme. La r\'eciproque est vraie si $I$ est nilpotent.
\end{proposition}

L'hypoth\`ese signifie que les homomorphismes
$$
I^n \otimes_A M \to I^n M
$$
sont des isomorphismes, d'o\`u aussit\^ot le fait que les
homomorphismes
$$
I^n / I^{n+1} \otimes_A M \to I^n M / I^{n+1} M
$$
sont des isomorphismes. R\'eciproquement, supposons cette condition
v\'erifi\'ee et $I$ nilpotent, prouvons $\Tor_1^A(M,A/I^n) =0$
pour tout $n$. C'est vrai pour $n$ grand, proc\'edons par
r\'ecurrence descendante sur~$n$, en le supposant prouv\'e pour
$n+1$. On a un diagramme commutatif
$$
\xymatrix{ &{M \otimes I^{n+1}}\ar[r]\ar[d]&{M \otimes I^n}\ar[r]\ar[d]&{M \otimes(I^n/I^{n+1})}\ar[r]\ar[d]&{0}\\{0}\ar[r]&{MI^{n+1}}\ar[r]&{MI^n}\ar[r]&{MI^n/MI^{n+1}}\ar[r]&{0}}
$$
o\`u
les lignes sont exactes. Par hypoth\`ese, la derni\`ere fl\`eche
verticale est un isomorphisme, et l'hypoth\`ese de r\'ecurrence
signifie aussi que la premi\`ere fl\`eche verticale l'est. Il en
est donc de m\^eme de la fl\`eche verticale m\'ediane, ce qui
ach\`eve la d\'emonstration.

La proposition suivante a \'et\'e d\'egag\'ee au moment du
S\'eminaire par Serre; elle permet des simplifications
substantielles dans le pr\'esent num\'ero.
\begin{proposition}
\label{IV.5.2}
Soient $A \to B$ un homomorphisme d'anneaux, $M$ un $A$-module. Les
conditions suivantes sont \'equivalentes:
\begin{enumerate}
\item[(i)] Pour tout $B$-module $N$, on a $\Tor_1^A(M,N) = 0$;
\item[(ii)] $\Tor_1^A(M,B) = 0$, et $M_{(B)} = M \otimes_A B$ est
$B$-plat.
\end{enumerate}
\end{proposition}

On a un isomorphisme fonctoriel
$$
M \otimes_A N = (M \otimes_A B) \otimes_B N
$$
qui exprime le premier membre, consid\'er\'e comme foncteur en
$M$, comme un compos\'e de deux foncteurs $M \mto M \otimes_A B$
et $P \mto P \otimes_B N$. Comme le premier transforme modules
libres sur~$A$ en modules libres sur~$B$, donc projectifs en
projectifs, on a la suite spectrale des foncteurs compos\'es
$$
\Tor_n^A(M,N) \From \Tor_p^B(\Tor_q^A(M,B),N)
$$
d'o\`u une suite exacte pour les termes de bas degr\'e
$$
0 \from \Tor_1^B(M \otimes_A B,N) \from \Tor_1^A(M,N)
\from \Tor_1^A(M,B) \otimes_A N
$$
Si (i) est v\'erifi\'e, alors on conclut de cette suite exacte
$\Tor_1^B(M \otimes_A B,N) = 0$ pour tout $N$, \ie $M \otimes_A B$
est $B$-plat, d'o\`u (ii). Si inversement (ii) est v\'erifi\'e,
alors dans la suite exacte les termes entourant $\Tor_1^A(M,N)$ sont
nuls, donc on a (i).

\begin{corollaire}
\label{IV.5.3}
Supposons que $B = A/I$, alors les conditions pr\'ec\'edentes
\'equivalent \`a la suivante:
\begin{enumerate}
\item[(iii)] $\Tor_1^A(M,N) = 0$ pour tout $A$-module $N$ annul\'e
par une puissance de~$I$.
\end{enumerate}
\end{corollaire}

En effet, (i) signifie qu'il en est ainsi si $N$ est annul\'e par
$I$. On en d\'eduit (iii) en appliquant l'hypoth\`ese aux
$I^nN/I^{n+1}N$.

\begin{corollaire}
\label{IV.5.4}
Sous les conditions de~\ref{IV.5.3}, les conditions envisag\'ees
impliquent que
l'ho\-mo\-mor\-phisme fonctoriel
\begin{equation*}
\label{eq:IV.5.*}
\tag{$*$} {\gr_I^0(M) \otimes_{A/I} \gr_I(A) \to \gr_I(M)}
\end{equation*}
est un isomorphisme, et que $M \otimes_A A/I$ est plat sur~$A/I$.
\end{corollaire}

Il suffit d'appliquer (iii) et \eqref{IV.5.1}. Utilisant la
r\'eciproque de~\ref{IV.5.1} dans le cas $I$ nilpotent, on trouve:
\begin{corollaire}
\label{IV.5.5}
Soient $A$ un anneau muni d'un id\'eal nilpotent $I$, $M$ un
$A$-module. Les conditions suivantes sont \'equivalentes:
\begin{enumerate}
\item[(i)] $M$ est $A$-plat
\item[(ii)] $M \otimes_A A/I$ est $A/I$-plat, et $\Tor_1^A(M,A/I) = 0$
\item[(iii)] $M \otimes_A A/I$ est $A/I$-plat, et l'homomorphisme
canonique \eqref{eq:IV.5.*} sur les gradu\'es est un isomorphisme.
\end{enumerate}
\end{corollaire}

En effet, ce sont respectivement les conditions (iii) et (ii)
pr\'ec\'edentes, et celles de corollaire~\ref{IV.5.4}.

Ne supposons plus $I$ nilpotent, alors on aura seulement a priori
dans~\ref{IV.5.5} les implications (i)$\To$(ii)$\To$(iii). D'autre part, comme la condition (iii) reste
stable en divisant par une puissance de~$I$, on voit en vertu
de~\ref{IV.5.5} qu'elle implique
\begin{enumerate}
\item[(iv)] \emph{pour tout entier $n$, $M \otimes A/I^n$ est plat sur~$A/I^n$.}
\end{enumerate}
On se propose de donner des conditions moyennant lesquelles on peut
en conclure~(i), \ie que $M$ est $A$-plat. Je dis qu'il suffit pour
ceci que $A$ soit noeth\'erien et que~$M$ satisfasse la condition de
finitude suivante: \emph{pour tout module de type fini $N$ sur~$A$, $M
\otimes_A N$ est séparé pour la topologie
$I$-préadique}. (Il suffirait de le v\'erifier si~$N$ est un
id\'eal de type fini dans $A$). En effet, prouvons que sous ces
conditions, si $N' \to N$ est un monomorphisme de modules de type
fini, $M \otimes_A N' \to M \otimes_A N$ est un monomorphisme. Il
suffit en effet de montrer que le noyau est contenu dans les $I^n(M
\otimes_A N') = \Im(M \otimes_A I^n N' \to M \otimes_A N')$, ou encore
dans les $\Im(M \otimes_A V'_n \to M \otimes_A N') = \Ker (M \otimes_A
N' \to M \otimes_A (N'/V'_n))$, o\`u $V'_n$ parcourt un syst\`eme
fondamental d\'enombrable de voisinages de 0 dans $N'$ (muni de sa
topologie $I$-adique). D'apr\`es le th\'eor\`eme de Krull, la
topologie $I$-adique de~$N'$ est induite par celle de~$N$, on peut
donc prendre $V'_n = N' \cap I^n N$. Consid\'erons alors le
diagramme commutatif
$$
\xymatrix{{M \otimes_A N'}\ar[r]\ar[d]&{M \otimes_A (N'/V'_n)}\ar[d]\\{M \otimes_A N}\ar[r]&{M \otimes_A (N / I^n N)}}
$$
Comme $N'/V'_n$ et $N/I^n N$ sont annul\'es par $I^n$, le
deuxi\`eme homomorphisme vertical s'identifie \`a celui d\'eduit
de l'homomorphisme \emph{injectif} $N'/V'_n \to N/I^n N$ en
tensorisant sur~$A/I^n$ avec le $(A/I^n)$-module \emph{plat} $M
\otimes_A A/I^n$, il est donc \emph{injectif}. Par suite, le noyau de
$M \otimes_A N' \to M \otimes_A N$ est contenu dans celui de~$M
\otimes_A N' \to M \otimes_A (N'/V'_n),$ ce qu'on voulait.

La condition de \og finitude\fg envisag\'ee sur~$M$ est
v\'erifi\'ee en particulier si $M$ est un module de type fini sur
une $A$-alg\`ebre
noeth\'erienne
$B$ telle que $IB$ soit contenu
dans le radical de~$B$: en effet, alors $M \otimes_A N$ est un module
de type fini sur~$B$ pour tout module de type fini $N$ sur~$A$, donc
s\'epar\'e par Krull pour la topologie $I$-adique $=$ sa topologie
$(IB)$-adique. On trouve ainsi:

\begin{theoreme}
\label{IV.5.6}
Soient $A\to B$ un homomorphisme d'anneaux noeth\'eriens, $I$ un
id\'eal de~$A$ tel que $IB$ soit contenu dans le radical de~$B$, $M$
un $B$-module de type fini. Les conditions suivantes sont
\'equivalentes:
\begin{enumerate}
\item[(i)] $M$ est $A$-plat
\item[(ii)] $M\otimes_A A/I$ est $A/I$-plat, et $\Tor_1^A(M,A/I)=0$
\item[(iii)] $M\otimes_A A/I$ est $A/I$-plat, et l'homomorphisme
canonique
$$
\gr_I^0(M)\otimes_{A/I}\gr_I(A)\to\gr_I(M)
$$
est un isomorphisme.
\item[(iv)] Pour tout entier $n$, $M\otimes_A A/I^n$ est plat
sur~$A/I^n$.
\end{enumerate}
\end{theoreme}

Ce r\'esultat s'applique surtout lorsque $A$, $B$ sont des anneaux
\emph{locaux} noeth\'eriens, $A\to B$ un homomorphisme local, et~$I$
un id\'eal de~$A$ contenu dans son id\'eal maximal (et on peut
r\'eduire aussit\^ot~\ref{IV.5.6} \`a ce cas). Un cas
int\'eressant est celui o\`u $A/I$ est un corps, \ie $I$ maximal,
auquel cas la condition que $M\otimes_A(A/I)$ est plat sur~$A/I$
devient inutile; de plus, comme alors les $A/I^n$ sont des anneaux
locaux artiniens, la condition (iv) signifie que les
$M\otimes_A(A/I^n)$ sont \emph{libres} sur les~$A/I^n$.

\begin{corollaire}
\label{IV.5.7}
Soit
$A\to B$ un homomorphisme local d'anneaux locaux noeth\'eriens,
$u\colon M'\to M$ un homomorphisme de~$B$-modules de type fini,
supposons $M$ plat sur~$A$. Alors les conditions suivantes sont
\'equivalentes:
\begin{enumerate}
\item[(i)] $u$ est injectif, et $\Coker u$ est plat sur~$A$.
\item[(ii)] $u\otimes_A k\colon M'\otimes_A k\to M\otimes_A k$ est
injectif
\end{enumerate}
(o\`u $k$ d\'esigne le corps r\'esiduel de~$A$).
\end{corollaire}

(i)$\To$(ii) en vertu de~\ref{IV.1.1}, prouvons la
r\'eciproque. Tout d'abord $u$ est injectif, car il suffit de le
v\'erifier sur les gradues associ\'es, o\`u cela r\'esulte
d'un carr\'e commutatif que le lecteur \'ecrira. Soit $M''$ son
$\Coker$, on a donc une suite exacte
$$
0\to M'\to M\to M''\to 0
$$
d'o\`u par la suite exacte des $\Tor$, compte tenu de
l'hypoth\`ese~(ii) et de~$\Tor_1^A(M,k)=0$, la relation
$\Tor_1^A(M'',k)=0$, donc $M''$ est plat sur~$A$ par le
th\'eor\`eme~\ref{IV.5.6}.

\begin{corollaire}
\label{IV.5.8}
Sous les conditions de~\ref{IV.5.6}, soit $J$ un id\'eal de~$B$
contenant $IB$ et contenu dans le radical. Soient $\hat A$ le
compl\'et\'e $I$-adiques de~$A$ et~$\hat B$ et~$\hat M$ les
compl\'et\'es $J$-adiques de~$B$ et~$M$. Pour que $M$ soit
$A$-plat, il faut et il suffit que $M$ soit $\hat A$-plat.
\end{corollaire}

(N.B. la suffisance r\'esulterait d\'ej\`a facilement
de~\ref{IV.3.2}). On utilise le crit\`ere (iii) de~\ref{IV.5.6}
dans la situation $(A,B,I,M)$ et dans la situation $(\hat A,\hat
B,I\hat A, \hat M)$. On constate que les conditions obtenues pour
l'un et l'autre cas sont \'equivalentes, gr\^ace
\`a~\ref{IV.3.2}.

\begin{corollaire}
\label{IV.5.9}
Soient $A\to B\to C$ des homomorphismes locaux d'anneaux locaux
noe\-th\'e\-riens, $M$ un $C$-module de type fini (N.B. $C$
n'intervient que pour pouvoir mettre une condition de finitude sur~$M$). On suppose $B$ plat sur~$A$. Soit $k$ le corps r\'esiduel
de~$A$. Conditions \'equivalentes:
\begin{enumerate}
\item[(i)] $M$ est plat sur~$B$.
\item[(ii)] $M$ est plat sur~$A$, et $M\otimes_A k$ est plat sur~$B\otimes_A k$.
\end{enumerate}
\end{corollaire}

L'implication (i)$\To$(ii) est triviale, prouvons
(ii)$\To$(i). On applique le crit\`ere~(iii)
de~\ref{IV.5.6} \`a $(B,C,\goth{m}B=I,M)$, comme
$M\otimes_B(B/I)=M\otimes_B(B\otimes_A k) = M\otimes_A k$, la
premi\`ere condition de ce crit\`ere signifie pr\'ecis\'ement
que $M\otimes_A k$ est plat sur~$B\otimes_A k$,
parfait. La deuxi\`eme condition du crit\`ere est
v\'erifi\'ee parce que $M$ est plat sur~$A$ et~$B$ plat sur~$A$,
par une formule d'associativit\'e du produit tensoriel. --- Bien
entendu, se r\'ef\'erant \`a~\ref{IV.5.5} au lieu
de~\ref{IV.5.6}, on obtient un \'enonc\'e analogue sans condition
noeth\'erienne et de finitude, quand on suppose en revanche que
l'id\'eal~$\goth{m}$ de~$A$ est nilpotent. (Le fait que $\goth{m}$
ait \'et\'e pris maximal n'est d'ailleurs pas intervenu; mais
c'est en un sens le cas \og $\goth{m}$ maximal\fg qui est \og le meilleur
possible\fg).

\section{Morphismes plats et ensembles ouverts}
\label{IV.6}

Rappelons d'abord quelques r\'esultats sur les ensembles
constructibles, qui sont d'ailleurs d\'emontr\'es dans des notes
en circulation du S\'eminaire Dieudonn\'e-Rosenlicht sur les
Sch\'emas\footnote{\Cf EGA~0$_\mathrm{III}$~9, EGA~IV~1.8 et~1.10.}.

Soit $X$ un espace topologique. On dit avec Chevalley qu'une partie
de~$X$ est \emph{constructible}
\index{constructible (partie)|hyperpage}si elle est r\'eunion finie de parties localement ferm\'ees.

\begin{lemme}
\label{IV.6.1}
Soit $X$ un espace topologique noeth\'erien, soit $Z$ une partie
de~$X$. Pour que~$Z$ soit constructible, il faut et il suffit que
pour toute partie ferm\'ee irr\'eductible $Y$ de~$X$, $Z\cap Y$
est non dense dans $Y$ ou contient une partie ouverte non vide de
l'espace~$Y$.
\end{lemme}

On en d\'eduit, utilisant un lemme bien connu d'Alg\`ebre
Commutative:

\begin{lemme}[Chevalley]
\label{IV.6.2}
Soit $f\colon X\to Y$ un morphisme de type fini de
pr\'esch\'emas, avec $Y$ noeth\'erien. Alors $f(X)$ est
constructible.
\end{lemme}

\begin{lemme}
\label{IV.6.3}
Soient $X$ un espace topologique noeth\'erien dont toute partie
ferm\'ee irr\'eductible admet un point g\'en\'erique, $U$ une
partie constructible de~$X$, $x\in X$. Pour que $U$ soit un voisinage
de~$x$, il faut et il suffit que toute g\'en\'erisation $y$ de~$x$
(\ie tout $y\in X$ tel que $x\in \bar{y}$) soit dans~$U$.
\end{lemme}

En particulier

\begin{corollaire}
\label{IV.6.4}
Soit $X$ un espace topologique noeth\'erien dont toute partie
ferm\'ee irr\'eductible admet un point g\'en\'erique, $U$ une
partie de~$X$. Pour que $U$ soit ouverte, il faut et il suffit
qu'elle satisfasse les deux conditions suivantes:
\begin{enumerate}
\item[(a)] $U$ contient toute g\'en\'erisation de chacun de ses
points
\item[(b)] si $x\in U$, alors
$U\cap\bar{x}$ contient une partie ouverte non vide de
l'espace~$\bar{x}$.
\end{enumerate}
\end{corollaire}

En effet, $U$ est n\'ecessairement constructible gr\^ace
\`a~\ref{IV.6.1}, et on applique le crit\`ere~\ref{IV.6.2} qui
prouve que $U$ est un voisinage de chacun de ses points.

\begin{corollaire}
\label{IV.6.5}
Soit $f\colon X\to Y$ un morphisme de type fini de pr\'esch\'emas,
avec $Y$ localement noeth\'erien, $x$ un point de~$X$, $y=f(x)$.
Pour que $f$ transforme toute voisinage de~$x$ en un voisinage de~$y$,
il faut et il suffit que pour toute g\'en\'erisation $y'$ de~$y$,
il existe une g\'en\'erisation $x'$ de~$x$ telle que $f(x')=y'$.
\end{corollaire}

On peut \'evidemment supposer que $X$ et~$Y$ sont affines, donc
noeth\'eriens. La condition est suffisante, car il suffit de
prouver que $f(X)$ est un voisinage de~$y$, or $f(X)$ est
constructible par~\ref{IV.6.1}, et il suffit d'appliquer le
crit\`ere~\ref{IV.6.3}. La condition est n\'ecessaire, car soit
$Y'=\overline{y'}$, et soit $F$ la r\'eunion des composantes
irr\'eductibles de~$f^{-1}(Y')$ qui ne contiennent pas $x$. Alors
$X-F$ est un voisinage ouvert de~$x$, donc son image est un voisinage
de~$y$, et a fortiori contient $y'$, donc il existe $x'_1\in X-F$ tel
que $f(x'_1)=y'$. Consid\'erons une composante irr\'eductible
de~$f^{-1}(Y')$ contenant $x'_1$, elle contient n\'ecessairement $x$
(car autrement elle serait contenue dans $F$), soit $x'$ son point
g\'en\'erique. C'est une g\'en\'erisation de~$x$, et $f(x')$
est une g\'en\'erisation de~$f(x'_1)=y'$ contenue dans~$Y'$, donc
est \'egal \`a $y'$, cqfd.

\begin{theoreme}
\label{IV.6.6}
Soient $f\colon X\to Y$ un morphisme localement de type fini, avec~$Y$
localement noeth\'erien, $F$ un faisceau coh\'erent sur~$X$ de
support~$X$, plat par rapport \`a~$Y$. Alors $f$ est un morphisme
ouvert (\iev transforme ouverts en ouverts).
\end{theoreme}

Il suffit de prouver le crit\`ere~\ref{IV.6.5} pour tout point $x\in
X$. Or les g\'en\'erisations $x'$ de~$x$ correspondent aux
id\'eaux premiers de~$\cal{O}_x$, celles~$y'$ de~$y$ correspondent
aux id\'eaux premiers de~$\cal{O}_y$, et il faut donc v\'erifier
que tout id\'eal premier de~$\cal{O}_y$ est induit par une id\'eal
premier de~$\cal{O}_x$. Or $F_x$ est un $\cal{O}_x$-module non nul et
$\cal{O}_y$ plat, donc fid\`element plat sur~$\cal{O}_y$
par~\ref{IV.2.2}. On peut donc appliquer~\ref{IV.2.3}, ce qui
ach\`eve la d\'emonstration.

\begin{remarquesstar}
Comme la platitude se conserve par extension de la base, on voit que
sous les conditions de~\ref{IV.6.5} $f$ est m\^eme
\emph{universellement ouvert}. J'ignore cependant, lorsque $Y$ est
int\`egre et $X$ de type fini sur~$Y$, si $f$ induit sur toute
composante $X_i$ de~$X$ un morphisme ouvert, ou m\^eme seulement
\'equidimensionnel\footnote{La r\'eponse \`a la deuxi\`eme question est
affirmative, celle \`a la premi\`ere n\'egative m\^eme si $f$
est \'etale; \cf EGA~IV~12.1.1.5 et EGA~Err$_{\rm IV}$~33.},
\ie
dont toutes les composantes des fibres ont m\^eme dimension (on sait
seulement que $X_i$ \emph{domine} $Y$). La question est li\'ee
\`a la suivante: soit $A\to B$ un homomorphisme local d'anneaux
locaux noeth\'eriens, tel que $B$ soit plat sur~$A$ et~$\goth{m}B$
soit un id\'eal de d\'efinition de~$B$, (ce qui implique
d'ailleurs $\dim B=\dim A$). Est-il vrai que pour tout id\'eal
premier minimal $\goth{p}_i$ de~$B$, on a $\dim B/\goth{p}_i=\dim B$?
Signalons seulement que la r\'eponse \`a la premi\`ere question
est n\'egative quand on remplace l'hypoth\`ese de platitude
de~\ref{IV.6.5} par la seule hypoth\`ese que $f$ soit
universellement ouvert.
\end{remarquesstar}

\begin{lemme}
\label{IV.6.7}
Soient $A$ un anneau int\`egre noeth\'erien, $B$ une
$A$-alg\`ebre de type fini, $M$ un $B$-module de type fini. Alors
il existe un \'el\'ement non nul $f$ de~$A$ tel que $M_f$ soit un
module libre (a fortiori plat) sur~$A_f$.
\end{lemme}

Soit $K$ le corps des fractions de~$A$, alors $B\otimes_A K$ est une
alg\`ebre de type fini sur~$K$, et $M\otimes_A K$ un module de type
fini sur cette derni\`ere. Soit $n$ la dimension du support de ce
module, nous raisonnerons par r\'ecurrence sur~$n$.
Si $n<0$,
\ie si $M\otimes_A K=0$, alors prenant un nombre fini de
g\'en\'erateurs de~$M$ sur~$B$, on voit qu'il existe un $f\in A$
qui annule ces g\'en\'erateurs, donc $M$, d'o\`u $M_f=0$ et on a
gagn\'e. Supposons $n\ge 0$. On sait que le $B$-module~$M$ admet
une suite de composition dont les quotients successifs sont isomorphes
\`a des modules $B/\goth{p}_i$, les $\goth{p}_i$ \'etant des
id\'eaux premiers de~$B$. Comme une extension de modules libres est
libre, on est ramen\'e au cas o\`u $M$ lui-m\^eme est de la
forme $B/\goth{p}$, ou encore identique \`a $B$, $B$ \'etant une
$A$-alg\`ebre \emph{intègre}. Appliquant le lemme de
normalisation de Noether \`a la $K$-alg\`ebre $B\otimes_A K$, on
voit facilement qu'il existe un \'el\'ement $f$ non nul de~$A$ tel
que $B_f$ soit entier sur le sous-anneau $A_f[t_1,\dots,t_n]$, o\`u
les $t_i$ sont des ind\'etermin\'ees. Donc on peut d\'ej\`a
supposer $B$ est entier sur~$C=A[t_1,\dots,t_n]$, c'est donc un
$C$-module de type fini sans torsion. Soit $m$ son rang, il existe
donc une suite exacte de~$C$-modules:
$$
0\to C^m\to B\to M'\to 0
$$
o\`u $M'$ est un $C$-module de torsion. Il s'ensuit que la
dimension de Krull du~$C\otimes_A K$-module $M'\otimes_A K$ est
strictement inf\'erieure \`a celle~$n$ de~$C\otimes_A K$.
D'apr\`es l'hypoth\`ese de r\'ecurrence, il s'ensuit que, \`a
condition de localiser par rapport \`a un $f$ non nul convenable
de~$A$, on peut supposer que $M'$ est un $A$-module
libre. D'autre part, $C^m$ est un $A$-module libre. Donc $B$ est
alors un $A$-module libre, on a fini.

\begin{lemme}
\label{IV.6.8}
Soient $A$ un anneau noeth\'erien, $B$ une alg\`ebre de type fini
sur~$A$, $M$~un $B$-module de type fini, $\goth{p}$ un id\'eal
premier de~$B$, $\goth{q}$ l'id\'eal premier qu'il induit sur~$A$. On suppose $M_\goth{p}$ plat sur~$A_{\goth{q}}$ (ou sur~$A$,
c'est pareil). Alors il existe un $g\in B-\goth{p}$ tel que
\begin{enumerate}
\item[(a)] $(M/\goth{q}M)_g$ est plat sur~$A/\goth{q}$.
\item[(b)] $\Tor_1^A(M,A/\goth{q})_g=0$.
\end{enumerate}
\end{lemme}
En effet, appliquant~\ref{IV.6.7} \`a $(A/\goth(q), B/\goth{q}B,
M/\goth{q}M)$ on voit d'abord qu'il existe un~$f$ dans $A-\goth{q}$
tel que $(M/\goth{q}M)_f$ soit plat sur~$A/\goth{q}$. D'autre part,
comme $M_\goth{p}$ est plat sur~$A$, on a
$\Tor_1^A(M,A/\goth{q})_\goth{p}=\Tor_1^A(M_\goth{p},A/\goth{q})=0$,
donc comme $\Tor_1^A(M,A/\goth{q})$ est un $B$-module de type fini, il
existe un $g\in B-\goth{p}$ tel qu'on ait (b). On peut alors
(rempla\c cant $g$ par $gf$) supposer qu'on a en m\^eme temps (a),
ce qui prouve le corollaire.

\begin{corollaire}
\label{IV.6.9}
Avec les notations de~\ref{IV.6.8}, pour tout id\'eal premier
$\goth{p}'$ de~$B$ contenant~$\goth{p}$ et ne contenant pas~$g$,
$M_{\goth{p}'}$ est plat sur~$A$ (ou, ce qui revient au m\^eme, sur~$A_{\goth{q}'}$, o\`u $\goth{q}'$ est l'id\'eal premier de~$A$
induit par~$\goth{p}')$.
\end{corollaire}
Il suffit d'appliquer le crit\`ere~\ref{IV.5.6} (ii) au syst\`eme
$(A, B_{\goth{q}'}, \goth{q}, M_{\goth{q'}})$, en utilisant la
localisation des~$\Tor$.
\begin{theoreme}
\label{IV.6.10}
Soit $f\colon X\to Y$ un morphisme de type fini, avec $Y$ localement
noeth\'erien, et soit $F$ un faisceau coh\'erent sur~$X$. Soit $U$
l'ensemble des points $x\in X$ tels que $F_x$ soit plat sur~$\cal{O}_{f(x)}$. Alors $U$ est un ensemble \emph{ouvert}.
\end{theoreme}

\subsubsection*{D\'emonstration} On peut supposer $X$ et $Y$ affines, d'anneaux $B$
et $A$, donc $F$ d\'efini par un $B$-module $M$ de type fini. On
applique le crit\`ere~\ref{IV.6.4}. La condition (a) est
v\'erifi\'ee trivialement par \ref{IV.1.2}~(i), reste \`a
v\'erifier la condition (b) de~\ref{IV.6.4}. C'est ce qui a
\'et\'e fait dans le lemme~\ref{IV.6.8} et
corollaire~\ref{IV.6.9}.

Dans beaucoup de questions, la forme plus faible suivante du
th\'eor\`eme~\ref{IV.6.10} est suffisante (qui r\'esulte
d\'ej\`a du lemme~\ref{IV.6.7}, et ne n\'ecessite donc ni la
technique
des constructibles, ni le th\'eor\`eme~\ref{IV.5.6}):
\begin{corollaire}
\label{IV.6.11}
Sous les conditions de~\ref{IV.6.10}, si on suppose $Y$ int\`egre,
alors il existe un ouvert non vide~$V$ dans $Y$ tel que $F$ soit plat
relativement \`a $Y$ en tous les points de~$f^{-1}(V)$.
\end{corollaire}
En effet, l'ensemble ouvert $U$ contient la fibre du point
g\'en\'erique de~$Y$ (puisque l'anneau local de ce point est un
corps), donc il contient un ouvert de la forme $f^{-1}(V)$, $X$
\'etant de type fini sur~$Y$. De~\ref{IV.6.11}, on conclut aussi
facilement le r\'esultat suivant, o\`u~$Y$ est suppos\'e
noeth\'erien (mais pas n\'ecessairement int\`egre): il existe
une partition de~$Y$ en des parties localement ferm\'ees $Y_i$
telles que (munissant $Y_i$ de la structure r\'eduite induite) $F$
induise sur chaque $X_i=X\times_Y Y_i$ un faisceau plat par rapport
\`a~$Y_i$.

